% Sample-construction tables for the response to reviewers.
% Styled to match templateV6.1.tex (ametsocV6.1 with makecell, xcolor, colortbl, array, booktabs).
% Needs \definecolor{rowgray}{RGB}{235,235,235}, which the manuscript already defines.
% Counts reproduced 2026-10-07 from the archived data in /glade/campaign/cisl/aiml/cat_flood.
%   Tables 1-3: intermediate_data/chunked_data_2023_2024, CAT_home_loc_20251010,
%               home_loc_by_row_group_counties.parquet, analysis_ready_data_20251121[_mrms].
%   Table 4:    jcelis pipeline (04_california_residents[-2023].py, Paired_Analysis.ipynb,
%               Common_UIDs.ipynb, POI_Fuel.py) and
%               ancillary_data/home_locations/2024/TAO_filtered_home_loc_by_row_group.parquet.

\begin{table*}[ht]
\centering
\small
\caption{Sample construction for the five-Saturday paired analysis, from the full device universe to the analytic sample, by county of inferred home location.}
\label{tab:sample_construction}
\begin{tabular}{|c|l|r|r|r|r|}
\hline
\rowcolor{rowgray}
\textbf{Step} & \textbf{Criterion} & \textbf{Los Angeles} & \textbf{Ventura} & \textbf{Santa Barbara} & \textbf{Total} \\
\hline
0 & Devices in data extract (92 days)$^{a}$ & \multicolumn{3}{c|}{---} & 75,017,659 \\
1 & Home location inferred$^{b}$ & \multicolumn{3}{c|}{---} & 42,039,428 \\
2 & Home ZIP code in study counties$^{c}$ & 10,255,396 & 924,575 & 329,041 & 11,509,012 \\
3 & $\geq$1 ping on 30 Mar 2024 (event Saturday) & 1,083,719 & 111,108 & 38,628 & 1,233,455 \\
4 & $\geq$1 stay location on 30 Mar 2024$^{d}$ & 575,453 & 60,996 & 21,125 & 657,574 \\
5 & Steps 2 and 4 met on all five Saturdays$^{e}$ & \textbf{71,782} & \textbf{8,679} & \textbf{3,042} & \textbf{83,503} \\
\hline
\rowcolor{rowgray}
\multicolumn{6}{|l|}{\textit{Exposure groups within the analytic sample (step 5)}} \\
\hline
 & Home inside a Flash Flood Warning polygon & 31,091 & 157 & 922 & 32,170 \\
 & MRMS 24-h precipitation at home $\geq$25 mm & 69,206 & 6,908 & 1,706 & 77,820 \\
 & MRMS 24-h precipitation at home $\geq$38 mm & 26,104 & 816 & 609 & 27,529 \\
 & MRMS 24-h precipitation at home $\geq$50 mm & 592 & 156 & 242 & 990 \\
\hline
 & Retention, step 5 / step 4 (\%) & 12.5 & 14.2 & 14.4 & 12.7 \\
\hline
\end{tabular}

\vspace{4pt}
\parbox{0.95\textwidth}{\footnotesize
$^{a}$Unique advertising IDs in the Veraset extract covering 15 Mar--29 Apr of 2023 and 2024 (12,207,313,402 ping records). No accuracy, speed, or ping-count filters were applied to the pings.
$^{b}$Stay locations ($\geq$10 min within 200 m; gaps $>$24 h end a stay) clustered at 100 m. Home is the cluster with the largest cumulative dwell time across all 92 days pooled. Devices with no stay location get no home.
$^{c}$345 ZIP codes assigned to Los Angeles, Ventura, and Santa Barbara counties. The eight ZIP codes listed under two counties are assigned to Ventura County.
$^{d}$At least one stay of $\geq$20 min within 200 m on that day, needed to compute daily mobility metrics.
$^{e}$The device meets step 4 on each of 16 Mar, 23 Mar, 30 Mar, 6 Apr, and 13 Apr 2024. The union of devices meeting step 4 on any Saturday is 1,882,357. Of these, 911,220, 483,678, 264,902, 139,054, and 83,503 meet it on exactly 1, 2, 3, 4, and 5 Saturdays.}
\end{table*}

\begin{table*}[ht]
\centering
\small
\caption{Daily coverage of tri-county residents (Table~\ref{tab:sample_construction}, step 2) on each analysis Saturday. Kept devices had at least one stay location that day (step 4). Device OS comes from the advertising-ID type (IDFA = iOS, AAID = Android). Fewer than 0.1\% of devices have another ID type.}
\label{tab:daily_coverage}
\begin{tabular}{|l|r|r|r|r|r|r|}
\hline
\rowcolor{rowgray}
 & \textbf{Devices with} & \textbf{Devices with} & & \multicolumn{2}{c|}{\textbf{iOS share (\%)}} & \textbf{Median pings} \\
\rowcolor{rowgray}
\textbf{Saturday} & \textbf{$\geq$1 ping} & \textbf{$\geq$1 stay} & \textbf{Kept (\%)} & \textbf{Kept} & \textbf{Dropped} & \textbf{kept / dropped} \\
\hline
16 Mar 2024 & 1,334,842 & 748,109 & 56.0 & 75.1 & 65.6 & 17 / 4 \\
23 Mar 2024 & 1,365,737 & 755,534 & 55.3 & 70.7 & 64.9 & 11 / 3 \\
\textbf{30 Mar 2024 (event)} & 1,233,455 & 657,574 & 53.3 & 70.8 & 64.0 & 12 / 3 \\
6 Apr 2024 & 1,253,517 & 695,797 & 55.5 & 62.8 & 54.4 & 14 / 3 \\
13 Apr 2024 & 1,386,172 & 789,999 & 57.0 & 52.6 & 49.8 & 13 / 3 \\
\hline
\end{tabular}
\end{table*}

\begin{table}[ht]
\centering
\small
\caption{Retained versus excluded devices among tri-county residents with a stay location on the event Saturday (30 Mar 2024; $N = 657{,}574$). Excluded devices had a stay location on 30 Mar but not on all four baseline Saturdays.}
\label{tab:retained_excluded}
\begin{tabular}{|l|r|r|}
\hline
\rowcolor{rowgray}
 & \textbf{Retained} & \textbf{Excluded} \\
\hline
Devices & 83,503 & 574,071 \\
\hline
\rowcolor{rowgray}
\multicolumn{3}{|l|}{\textit{Data density}} \\
\hline
Pings on 30 Mar, median (mean) & 22 (42.1) & 11 (19.6) \\
Stay locations on 30 Mar, median & 2 & 1 \\
Cumulative dwell at home cluster, median (days)$^{a}$ & 10.1 & 2.5 \\
iOS share (\%) & 74.6 & 70.3 \\
\hline
\rowcolor{rowgray}
\multicolumn{3}{|l|}{\textit{Geography (\% of devices)}} \\
\hline
Los Angeles / Ventura / Santa Barbara & 86.0 / 10.4 / 3.6 & 87.7 / 9.1 / 3.1 \\
\hline
\rowcolor{rowgray}
\multicolumn{3}{|l|}{\textit{Mobility}} \\
\hline
Straight-line distance on 30 Mar, median (km) & 9.3 & 1.6 \\
Mean baseline-Saturday distance, median (km)$^{b}$ & 24.6 & 5.7 \\
\hline
\rowcolor{rowgray}
\multicolumn{3}{|l|}{\textit{Exposure}} \\
\hline
Home inside a Flash Flood Warning polygon (\%) & 38.5 & 44.8 \\
MRMS 24-h precipitation at home, median (mm) & 36.4 & 36.4 \\
\hline
\end{tabular}

\vspace{4pt}
\parbox{0.95\columnwidth}{\footnotesize
$^{a}$Total time in the home cluster over the 92-day home-inference window, a proxy for how long the device was observed.
$^{b}$Mean over the baseline Saturdays (16 Mar, 23 Mar, 6 Apr, 13 Apr) on which the device has a stay location. Excluded: $n = 422{,}631$ devices with at least one baseline Saturday.}
\end{table}

% Retention across the 345 home ZIP codes (step 5 / step 4 on 30 Mar):
% median 12.6%, IQR 9.9--14.7%, range 0--27.3%.

\begin{table*}[ht]
\centering
\small
\caption{Sample construction for the residency-based paired analysis, the 2023--2024 overlap sample, and the point-of-interest (POI) analysis.}
\label{tab:sample_construction_other}
\begin{tabular}{|c|l|r|}
\hline
\rowcolor{rowgray}
\textbf{Step} & \textbf{Criterion} & \textbf{Devices} \\
\hline
\rowcolor{rowgray}
\multicolumn{3}{|l|}{\textit{A. 2024 residents (14 Mar--29 Apr 2024, 47 days)}} \\
\hline
A0 & Devices in 2024 data extract$^{a}$ & 55,212,368 \\
A1 & $\geq$1 ping on $\geq$45 of 47 days (no spatial restriction) & 48,584 \\
A2 & Device-days with daily distance $\leq$500 mi; has daily metrics & 48,582 \\
A3 & Home ZIP code in study-area list$^{b}$ & 10,849 \\
A4 & Of A3, also present in the 2023 extract (any ping)$^{c}$ & 7,178 \\
\hline
\rowcolor{rowgray}
\multicolumn{3}{|l|}{\textit{B. 2023 residents (14 Mar--29 Apr 2023, 47 days)}} \\
\hline
B0 & Devices in 2023 data extract & 30,387,102 \\
B1 & $\geq$1 ping on $\geq$45 of 47 days (no spatial restriction)$^{d}$ & 36,476 \\
\hline
\rowcolor{rowgray}
\multicolumn{3}{|l|}{\textit{C. Persistent in both years (alternative overlap definitions)}} \\
\hline
C1 & A1 $\cap$ B1 (statewide) & 4,668 \\
C2 & A3 $\cap$ B1 (study area) & 1,182 \\
\hline
\rowcolor{rowgray}
\multicolumn{3}{|l|}{\textit{D. Supermarket POI analysis (15 Mar--29 Apr 2024)}} \\
\hline
D0 & Devices in combined 2023--2024 extract (Table~\ref{tab:sample_construction}, step 0) & 75,017,659 \\
D1 & $\geq$1 stay ($\geq$5 min, $\leq$12 h, within 50 m) at one of 588 supermarket POIs in the study counties$^{e}$ & 139,041 \\
\hline
\end{tabular}

\vspace{4pt}
\parbox{0.95\textwidth}{\footnotesize
$^{a}$From a 2024-only chunking of the raw daily files (14 Mar--29 Apr 2024). The two archived 2024 chunkings contain 55,212,368 and 55,453,054 devices.
$^{b}$Home locations from a separate stay-location home-inference run on 2024 data only (archived as \texttt{TAO\_filtered\_home\_loc\_by\_row\_group.parquet}), not the run behind Table~\ref{tab:sample_construction}. They are matched against a list of 568 ZIP codes, a few of which appear to lie outside the three study counties.
$^{c}$Requires only that the device appear in the 2023 extract, not that it meet the 2023 residency criterion (B1).
$^{d}$The manuscript reports 36,472. The output file and run log both give 36,476.
$^{e}$Matched stays: 234,703; supermarkets with $\geq$1 stay: 583. By county: Los Angeles 122,188, Ventura 11,984, Santa Barbara 5,247 devices. Applies no residency or home-location filter. The original manuscript reported 102,245 devices (172,205 stays) because only device row groups 0--55,212 of 75,018 were processed. The full rerun reproduces all 172,205 of those stays.}
\end{table*}
